\section*{ЗАКЛЮЧЕНИЕ}
\phantomsection
\addcontentsline{toc}{section}{Заключение}

В ходе выполнения курсовой работы проведён численный анализ влияния дискретного микрорельефа суперэллиптической формы на трибологические характеристики радиального подшипника скольжения. Решены следующие задачи:

\begin{enumerate}
    \item Выполнен обзор теоретических основ гидродинамической теории смазки и современных работ, посвящённых лазерному текстурированию поверхностей трения. Показано, что основным механизмом положительного влияния микрорельефа является микрогидродинамический эффект, возникающий на кромках углублений за счёт формирования локальных клиновых зазоров.

    \item Изучена методика численного решения уравнения Рейнольдса для подшипника конечной длины методом конечных разностей с использованием итерационной процедуры Гаусса--Зейделя и верхней релаксацией. Расчёт проведён на сетке $500 \times 500$ узлов при критерии сходимости $\delta < 10^{-5}$.

    \item Рассчитаны поля давления и интегральные характеристики (нагрузочная способность $F$, коэффициент трения $\mu$, расход смазки $Q$) для гладкого подшипника и подшипника с суперэллиптическими углублениями (тип~7) в диапазоне эксцентриситетов $\varepsilon = 0{,}05 - 0{,}80$. Геометрия подшипника: $R = 50$~мм, $c = 70$~мкм, $L = 80$~мм.

    \item Проведён сравнительный анализ гладкого и текстурированного подшипников. При характерном рабочем эксцентриситете $\varepsilon = 0{,}6$ получены следующие результаты:
    \begin{itemize}
        \item нагрузочная способность возросла на $7{,}28$\%: с $F = 10685{,}6$~Н до $F = 11464{,}0$~Н;
        \item коэффициент трения уменьшился на $7{,}80$\%: с $\mu = 0{,}01507$ до $\mu = 0{,}01389$;
        \item расход смазки увеличился на $0{,}67$\%: с $Q = 0{,}6739$~л/с до $Q = 0{,}6785$~л/с.
    \end{itemize}
\end{enumerate}

Полученные результаты свидетельствуют о том, что нанесение суперэллиптических углублений с параметрами $a = 2{,}55$~мм, $b = 2{,}35$~мм, $h_p = 10$~мкм ($h_p/c = 0{,}1$) приводит к улучшению рабочих характеристик подшипника. Одновременное повышение нагрузочной способности и снижение коэффициента трения при незначительном росте расхода смазки указывает на эффективность данного типа микрорельефа.

Суперэллиптический профиль углублений, описываемый степенной функцией четвёртого порядка ($\Delta h = h_p \sqrt{1 - r^4}$), по резкости стенок занимает промежуточное положение между эллипсоидальным и цилиндрическим профилями. Более крутые стенки по сравнению с эллипсоидом обеспечивают усиленный микрогидродинамический эффект на кромках, при этом отсутствие разрыва производной (в отличие от цилиндрического профиля) способствует плавному течению смазки без резких скачков давления.

На основании проведённого анализа можно рекомендовать применение суперэллиптического микрорельефа для подшипников скольжения нефтегазового оборудования, работающего в условиях жидкостного трения. Наибольшая эффективность текстурирования достигается в тяжелонагруженных узлах ($\varepsilon > 0{,}5$), где прирост нагрузочной способности и снижение трения проявляются наиболее существенно.
